% !TeX program = lualatex
\documentclass[tikz,border=6pt]{standalone}

\usepackage{tikz}
\usepackage{amsmath}

\begin{document}
\begin{tikzpicture}[x=1cm,y=1cm, line cap=round, line join=round, font=\small]

% ============================================================
% A) THAM SỐ HÓA
% ============================================================
\def\xMax{8}      % biên phải của lưới (đơn vị)
\def\yMax{7}      % biên trên của lưới (đơn vị)

% Tọa độ các điểm (theo hình)
\def\Ax{1}\def\Ay{2}
\def\Bx{2}\def\By{4}
\def\Cx{2}\def\Cy{1}
\def\Dx{4}\def\Dy{3}
\def\Ex{6}\def\Ey{0}
\def\Fx{7}\def\Fy{5}
\def\Gx{6}\def\Gy{6}

% Độ dài phần trục vượt ra ngoài lưới để giống hình
\def\xLeftExt{1.2}
\def\xRightExt{0.8}
\def\yDownExt{1.2}
\def\yUpExt{0.4}

% ============================================================
% B) LƯỚI VÀ TRỤC
% ============================================================
% Lưới trong góc phần tư I
\draw[step=1, black!35] (0,0) grid (\xMax,\yMax);

% Trục tọa độ (dày)
\draw[line width=1.0pt] (0,0) -- (\xMax,0);
\draw[line width=1.0pt] (0,0) -- (0,\yMax);

% Nhãn số trên trục (như hình: 2,4,6 trên Oy; 5 trên Ox)
\node[left]  at (0,2) {2};
\node[left]  at (0,4) {4};
\node[left]  at (0,6) {6};
\node[below] at (5,0) {5};

% Chú thích đơn vị
\node[below] at (3.5,-0.65) {$1\;\text{đơn vị} = 50\;\text{mile}$};

% ============================================================
% C) CÁC ĐIỂM VÀ NHÃN
% ============================================================
\fill (\Ax,\Ay) circle (2.2pt);
\fill (\Bx,\By) circle (2.2pt);
\fill (\Cx,\Cy) circle (2.2pt);
\fill (\Dx,\Dy) circle (2.2pt);
\fill (\Ex,\Ey) circle (2.2pt);
\fill (\Fx,\Fy) circle (2.2pt);
\fill (\Gx,\Gy) circle (2.2pt);

\node[left]  at (\Ax,\Ay) {A};
\node[left]  at (\Bx,\By) {B};
\node[below] at (\Cx,\Cy) {C};
\node[right] at (\Dx,\Dy) {D};
\node[below] at (\Ex,\Ey) {E};
\node[right] at (\Fx,\Fy) {F};
\node[above] at (\Gx,\Gy) {G};

\end{tikzpicture}
\end{document}
